\documentclass{article}

\usepackage{listings}
\usepackage[english]{babel}

\usepackage[
letterpaper,
top=2cm,
bottom=2cm,
left=3cm,
right=3cm,
marginparwidth=1.75cm
]{geometry}

\usepackage{amsmath}
\usepackage{graphicx}
\usepackage[colorlinks=true, allcolors=blue]{hyperref}

\title{Linear Regression using Gradient Descent}
\author{Truong Quang Nam}

\begin{document}

\maketitle

\section{Introduction}

Linear Regression is a supervised machine learning algorithm used to predict continuous values. The model attempts to find the best linear relationship between input variable $x$ and output variable $y$.

The equation of linear regression is:

\[
y = w_1x + w_0
\]

Where:

\begin{itemize}
    \item $w_1$ is the weight (slope)
    \item $w_0$ is the bias (intercept)
\end{itemize}

To optimize the parameters, Gradient Descent is used to minimize the loss function.

\section{Implementation}

\subsection{Loss Function}

The loss function used in this implementation is Mean Squared Error (MSE):

\[
L(w_0,w_1,x,y) = \frac{1}{2}(w_1x + w_0 - y)^2
\]

The total cost function is:

\[
J(w_0,w_1) = \frac{1}{n}\sum_{i=1}^{n}L(w_0,w_1,x_i,y_i)
\]

\subsection{Read Dataset}

The dataset is loaded from a CSV file.

\begin{lstlisting}
def readcsv(f):
    x = []
    y = []

    with open(f, "r") as file:
        lines = file.readlines()[1:]

        for line in lines:
            xi, yi = line.strip().split(",")

            x.append(float(xi))
            y.append(float(yi))

    return x, y
\end{lstlisting}

\subsection{Loss Function Implementation}

\begin{lstlisting}
def L(w0, w1, x, y):
    return 0.5 * (w1 * x + w0 - y) ** 2

def J(w0, w1, x, y):
    total = 0
    n = len(x)

    for i in range(n):
        total += L(w0, w1, x[i], y[i])

    return total / n
\end{lstlisting}

\subsection{Gradient Calculation}

The derivative of the loss function with respect to $w_0$ and $w_1$ is implemented below.

\begin{lstlisting}
def dw0(w0, w1, x, y):
    total = 0
    n = len(x)

    for i in range(n):
        total += (w1 * x[i] + w0 - y[i])

    return total / n

def dw1(w0, w1, x, y):
    total = 0
    n = len(x)

    for i in range(n):
        total += x[i] * (w1 * x[i] + w0 - y[i])

    return total / n
\end{lstlisting}

\subsection{Gradient Descent Algorithm}

The weights are updated iteratively until the loss difference becomes smaller than threshold $t$.

\[
w = w - r \times \frac{dJ}{dw}
\]

Where:

\begin{itemize}
    \item $r$ is the learning rate
    \item $\frac{dJ}{dw}$ is the gradient
\end{itemize}

\begin{lstlisting}
def gradient_descent(w0, w1, x, y, r, t, max_epoch=100000):
    epoch = 0

    old_loss = float("inf")
    new_loss = J(w0, w1, x, y)

    while abs(old_loss - new_loss) > t and epoch < max_epoch:

        old_loss = new_loss

        d0 = dw0(w0, w1, x, y)
        d1 = dw1(w0, w1, x, y)

        w0 = w0 - r * d0
        w1 = w1 - r * d1

        new_loss = J(w0, w1, x, y)

        epoch += 1

        if epoch % 100 == 0:
            print(
                f"Epoch {epoch} "
                f"w0:{w0:.4f} "
                f"w1:{w1:.4f} "
                f"Loss:{new_loss:.6f}"
            )

    return w0, w1
\end{lstlisting}

\section{Result}

Initial parameters:

\begin{itemize}
    \item $w_0 = 0$
    \item $w_1 = 1$
    \item Learning rate $r = 0.0001$
\end{itemize}

Example training output:

\begin{table}[h!]
\centering
\caption{Training Process}
\begin{tabular}{|c|c|c|c|}
\hline
Epoch & $w_0$ & $w_1$ & Loss \\ \hline
100 & 0.1093 & 2.0649 & 254.832 \\ \hline
200 & 0.2007 & 2.0634 & 253.997 \\ \hline
300 & 0.2919 & 2.0619 & 253.164 \\ \hline
400 & 0.3830 & 2.0603 & 252.335 \\ \hline
500 & 0.4739 & 2.0588 & 251.508 \\ \hline
\dots & \dots & \dots & \dots \\ \hline
100000 & 44.610 & 1.323 & 12.745081 \\ \hline
\end{tabular}
\end{table}

Final parameters after optimization:

\[
w_0 = 44.610
\]

\[
w_1 = 1.323
\]

\section{Conclusion}

In this labwork, Linear Regression was implemented from scratch using Python. Gradient Descent was applied to optimize the weight and bias parameters by minimizing the Mean Squared Error loss function.

The result shows that Gradient Descent successfully reduces the loss value over time and finds a better linear model for prediction.

\end{document}